\section{Experimental Setup}
\label{sec:experiments}

\subsection{Data, model, and protocol}
All retained comparisons are conducted on \textbf{Cohere command-r-plus-08-2024} with a fixed per-method budget cap of $B = 6$ logical API calls per example.  The dataset is \texttt{openai/gsm8k}.  Evaluation conditions are matched across all compared candidate-producing methods: every method operates under the same budget cap, provider, and model, so accuracy differences are not confounded by unequal per-method compute.

\paragraph{Matched-budget definition.} A comparison is \emph{matched-budget} if and only if every compared candidate-producing method is subject to the same maximum logical API call count $B$ per example, with the same provider and model.  No method in the retained comparison receives extra calls or a different model.  FTA is then applied as a post-generation selection policy over the logged candidate answers and frontier metadata.
This contract supports clean method-to-method comparison, but it does not imply total compute parity with a single baseline run when multiple candidate-producing methods are executed together.

\subsection{Cost and resource usage}
\label{sec:cost}

Table~\ref{tab:cost_budget} reports per-method resource usage measured from logged Aggregate-720 runs on Cohere command-r-plus-08-2024. Reported metrics include average logical API calls, average input tokens, average output tokens, and average latency in seconds.  Total tokens (input + output) can be derived by summing the respective columns.  Budget parameter $B = 6$ is shown for completeness.

FTA adds negligible selection overhead once candidate answers and frontier
metadata are available, but the main deployment cost is candidate-pool construction.
If a deployment must generate frontier, L1, S1, and TALE outputs from scratch, total
wall-clock and API cost depends on which candidate producers are executed and may
exceed a single method's $B=6$ run. In other words, FTA is a zero-extra-call selector
after generation, not a claim of one-run compute equivalence to all baselines combined.
In the logged Aggregate-720 runs, the frontier
method has higher average latency than L1, S1, or TALE.
We report cost in logical calls, tokens, and latency; monetary cost depends on provider pricing and is not used as the primary comparison metric.  Dollar-normalized cost analysis, latency percentile distributions, and multi-provider token cost benchmarks are identified as future work (Section~\ref{sec:limitations}).

\subsection{Parameter grounding}
\label{sec:param_grounding}

Table~\ref{tab:param_sensitivity} (Appendix) documents the key parameters and design choices used in the two-gate FTA policy, including their values, roles, and the robustness evidence available.  We report parameter grounding and robustness checks rather than a full continuous hyperparameter sweep; a full threshold sweep is left for future work.

\subsection{Primary splits and disjointness}
The primary aggregate (Aggregate-720) is formed from three disjoint sources: seed~41 (300 examples), seed~61 (120 examples), and seed~71 (300 examples).  Cross-source overlap checks confirm zero shared example IDs and zero shared question hashes, ensuring that the aggregate evidence is drawn from non-overlapping subpopulations.

\subsection{Statistics}
We report paired bootstrap confidence intervals \citep{efron1993bootstrap} ($n = 5{,}000$ resamples) on Final-300 and source-stratified aggregate confidence intervals on Aggregate-720.  The bootstrap resample count of 5,000 is standard for stable CI estimation at the sample sizes used.  All confidence intervals are two-sided at the 95\% level.  When shown, the \emph{bootstrap positive-delta fraction} is the fraction of bootstrap resamples in which the FTA delta is positive; it is a descriptive bootstrap summary and should not be interpreted as a classical frequentist $p$-value.
